% ===================================================================
%  圖表第 16 號 — 大公司。--- 集市。玩具。
%
%  Traduction, faite en 2026, en cantonais ecrit d'aujourd'hui, en
%  caracteres traditionnels, sur l'ido de texto/io. Regles de la
%  colonne : voir 10-toubiu-01.tex. Une seule numerotation. Dernier
%  tableau du livret.
%
%  LES JOUETS DONNENT LE MOT QUI RESUME TOUTE LA COLONNE :
%
%    公仔 (83)  la poupee, la figurine  Pekin : 娃娃
%
%  « 公仔 » ne designe pas seulement la poupee : c'est le mot de tout
%  ce qui est figure — 公仔麵 les nouilles instantanees a l'image sur
%  le paquet, 公仔書 la bande dessinee, 公仔箱 le poste de
%  television. Un mot qui fait souche est un mot vivant, et c'est
%  toute la these de la colonne en un caractere double.
%
%  LA MENAGERIE DE CARTON EN AJOUTE DEUX :
%
%    馬騮 (10)   le singe     Pekin : 猴子
%    格力狗 (14) le levrier   Pekin : 灵缇  — « greyhound », entre par
%                                            les courses de Hong Kong
%    蝠鼠 (35)   la chauve-souris  Pekin : 蝙蝠
%
%  ET LE MAGASIN LUI-MEME :
%
%    夾萬 (67)   le coffre-fort   Pekin : 保险柜
%    五桶櫃 (68) la commode
%    銀包 (88)   le porte-monnaie Pekin : 钱包
%    結他 (72)   la guitare       Pekin : 吉他
%    沙律        la salade        Pekin : 沙拉
%
%  « 夾萬 » veut dire « qui serre dix mille » — dix mille pieces —
%  et n'a rien a voir avec l'assurance dont le nord a tire son mot.
%
%  LA VOITURE D'ENFANT (69) EST RENDUE PAR 嬰兒車 ET NON PAR LE
%  BB車 DE L'USAGE, qui s'ecrit avec deux lettres latines : la
%  colonne tient en caracteres, et un sigle romain y ferait tache
%  la ou tout le reste tient debout tout seul.
%
%  ECART DECLARE : le bloc t16-15-1 rend la lettre (d) des deux
%  hemispheres, que le fac-simile ido compose « d) » sans parenthese
%  ouvrante alors que la planche la grave. renvoji.py le passe, la
%  raison etant inscrite dans APARTA.
%
%  CE TABLEAU CLOT LA COLONNE. Seize planches, 672 blocs, et un seul
%  obstacle repete : le chinois place le determinant devant le
%  determine, l'ido l'inverse, et le renvoi suit son mot. Trente-cinq
%  divergences au debut, aux tableaux 3 a 6 ; zero au tableau 15,
%  parce qu'il enumere. Le remede n'a jamais change — nommer d'abord
%  ce que l'ido nomme d'abord, qualifier ensuite par une seconde
%  proposition.
% ===================================================================

%%K t16-apar-1 apar
\VUsaut{4.0mm}
\VUtitre{15.0pt}{60}{圖表第 16 號}
\VUsaut{2.0mm}
\VUpk{13.2pt}{40}{大公司。--- 集市。}
\VUpk{13.2pt}{40}{玩具。}
\VUsaut{3.4mm}

%%K t16-01-1 p
1. --- 新年就快到嘞。啲大公司已經擺好晒啲\VUgras{玩具} \textsuperscript{(1)} 同賀年禮物。

%%K t16-01-2 p
我求我阿爺帶我去集市睇下呢個靚展覽，佢就應承咗。

%%K t16-02-1 p
2. --- 我哋第一個停低嘅係啲靚兵器盾牌，上面掛住\VUgras{步槍、}\VUgras{軍帽、}
\VUgras{軍刀、}\VUgras{佩劍帶}同一對\VUgras{肩章；}又或者\VUgras{頭盔、}\VUgras{胸甲}
\textsuperscript{(3)} 同\VUgras{手槍} \textsuperscript{(4)}。呢啲嘢我覺得好睇過部
\VUgras{幻燈機} \textsuperscript{(2)} 好多。

%%K t16-tit-1 sub
\VUsaut{3.0mm}
\VUpk{10.2pt}{40}{好睇嘅嘢。}
\VUsaut{2.6mm}

%%K t16-03-1 p
3. --- 同一個部門度有個\VUgras{哨兵} \textsuperscript{(5)} 企喺個\VUgras{崗亭}度，
好似守住成個紙板動物園噉。入面幾多隻動物呀：企喺架上面嘅\VUgras{鸚鵡}
\textsuperscript{(6)}、鼻上生隻角嘅\VUgras{犀牛} \textsuperscript{(7)}、\VUgras{馴鹿}
\textsuperscript{(8)}、\VUgras{鴕鳥} \textsuperscript{(9)}、扭住粒合桃嘅\VUgras{馬騮}
\textsuperscript{(10)}、叼住隻雞乸走嘅\VUgras{狐狸} \textsuperscript{(11)}、擘大兩排大牙嘅
\VUgras{鱷魚} \textsuperscript{(12)}、\VUgras{駱駝} \textsuperscript{(13)}、身型好靚嘅
\VUgras{格力狗} \textsuperscript{(14)}、\VUgras{豹} \textsuperscript{(15)}；跟住有兩隻我唔識嘅，
人哋話係\VUgras{黃鼠狼} \textsuperscript{(16)} 同\VUgras{鬣狗} \textsuperscript{(17)}。嗰度仲有
\VUgras{金錢豹} \textsuperscript{(18)} 同\VUgras{狼} \textsuperscript{(21)}；就喺隔籬有啲得意嘅
\VUgras{老鼠仔} \textsuperscript{(19)} 同細到不得了嘅樹膠\VUgras{老鼠} \textsuperscript{(20)}。
最後係\VUgras{河馬} \textsuperscript{(22)} 同

%%K t16-03-1 p suite
生住個駝峰嘅\VUgras{單峰駝} \textsuperscript{(23)}，收埋成套。要唔係隔籬有個靚到極嘅
營地，擺滿\VUgras{鉛兵} \textsuperscript{(29)} 吸引咗我，我實喺嗰度企多幾個鐘。

%%K t16-tit-2 sub
\VUsaut{4.0mm}
\VUpk{10.2pt}{40}{兵同仗。}
\VUsaut{2.8mm}

%%K t16-04-1 p
4. --- 我阿爺係\VUgras{傷殘軍人} \textsuperscript{(24)}，佢因為受咗\VUgras{傷}要退伍，
仲攞咗個\VUgras{軍功章；}佢最鍾意同我講佢打過嗰啲\VUgras{仗。}佢好開心噉解釋
呢隊縮細咗嘅小軍隊點樣調動。有一班真兵嚟逛集市：一個\VUgras{工兵}
\textsuperscript{(25)}、一個戴住\VUgras{貝雷帽}嘅\VUgras{阿爾卑斯獵兵} \textsuperscript{(26)}、
一個步兵\VUgras{司務長} \textsuperscript{(27)}，同一個袖上釘住金\VUgras{線}嘅
\VUgras{炮兵中士} \textsuperscript{(28)}——佢哋都停低喺我哋隔籬聽。

%%K t16-05-1 p
5. --- 我阿爺見到有咁威水嘅聽眾，威到爆，即席就砌咗個完整嘅作戰計劃出嚟。

%%K t16-05-2 p
座有\VUgras{城牆}同\VUgras{護城河}嘅要塞，畀\VUgras{敵軍}圍住；敵軍\VUgras{炮轟}
咗好耐之後，就準備\VUgras{強攻。}不過\VUgras{守軍}捱唔住餓，就想\VUgras{殺出去}
攻\VUgras{圍城軍}個\VUgras{營。}圍城軍喺\VUgras{營地}周圍死頂，打退咗呢場
\VUgras{進攻；}\VUgras{打贏}咗之後，佢哋就放出\VUgras{騎兵}去追\VUgras{打輸}嘅
攻擊方。攻擊方一路\VUgras{退，}一路將\VUgras{戰場}鋪滿\VUgras{死人}同\VUgras{傷者。}
\VUgras{抬擔架嘅}就將啲傷者抬入\VUgras{救護站，}交畀\VUgras{外科醫生}醫。

%%K t16-05-3 p
就算有幾個\VUgras{營}排成\VUgras{方陣}英勇抵抗，都改變唔到\VUgras{全軍覆沒；}

%%K t16-05-3 p suite
剩返嘅軍隊都\VUgras{走}晒，丟低一大堆\VUgras{俘虜。}敵軍入城，就\VUgras{贏}到徹底。
呢場仗大概就係噉完；\VUgras{停戰}之後，就簽得成\VUgras{和約}嘞。

%%K t16-tit-3 sub
\VUsaut{4.4mm}
\VUpk{10.2pt}{40}{賀年禮物。}
\VUsaut{2.8mm}

%%K t16-06-1 p
6. --- 幾個後生兵一路聽老兵講得咁生鬼一路笑，仲追住問多啲。我自己就
繞住個檔口行，睇下嗰把靚\VUgras{弩} \textsuperscript{(32)}，同埋張\VUgras{弓}
\textsuperscript{(33)} 連支\VUgras{箭} \textsuperscript{(31)}。喺嗰度我又見到另一批動物：
兩隻\VUgras{會唱歌嘅雀} \textsuperscript{(34)}——一隻機械\VUgras{夜鶯，}一隻扯開喉嚨唱嘅
\VUgras{林鶯；}仲有一列古古怪怪嘅：\VUgras{蝠鼠} \textsuperscript{(35)}、\VUgras{蟒蛇}
\textsuperscript{(36)}、\VUgras{龜} \textsuperscript{(37)}、\VUgras{海豹} \textsuperscript{(38)}、
\VUgras{鯨魚} \textsuperscript{(39)}，同一隻樣衰衰嘅\VUgras{鯊魚} \textsuperscript{(40)}。

%%K t16-07-1 p
7. --- 我冇望好耐，因為我見到我發夢都想要嗰樣嘢：一個靚到極嘅
\VUgras{木偶戲台} \textsuperscript{(41)}，有一流嘅\VUgras{佈景}同一幅靚到迷人嘅紅
\VUgras{幕。}\VUgras{台}上兩個演員好似喺度做緊一齣好好睇嘅\VUgras{戲}嘅\VUgras{角色，}
因為啲紙板\VUgras{觀眾}——有啲坐喺包廂，有啲坐喺\VUgras{大椅}上面，有啲坐喺
\VUgras{樂隊指揮}後面嘅\VUgras{堂座}——都拍爛手掌。

%%K t16-07-2 p
可惜嗰個戲台貴到飛起。

%%K t16-08-1 p
8. --- 何況要揀玩具真係難，因為啲櫥窗堆滿各式各樣：\VUgras{丑角}
\textsuperscript{(42)}、\VUgras{木偶} \textsuperscript{(43)}、\VUgras{白面小丑} \textsuperscript{(44)}、
拍緊翼嘅\VUgras{貓頭鷹} \textsuperscript{(45)}、識吹口哨嘅\VUgras{烏鶇} \textsuperscript{(46)}、
識吠嘅\VUgras{貴婦狗、}\VUgras{留聲機} \textsuperscript{(47)}，同啲\VUgras{益智玩具。}其中一件
玩具好逗到我：

%%K t16-noto-1 noto
\VUnotes{143.2mm}{%
\textit{(*) 睇圖表第 6 號。}%
}

%%K t16-08-1 p suite
佢做嘅係一場\VUgras{海戰} \textsuperscript{(48)}：一隻\VUgras{船}喺一笪種滿\VUgras{椰子樹}
嘅陸地附近，畀啲\VUgras{海盜}襲擊。

%%K t16-09-1 p
9. --- 我好易就睇晒呢啲奇景，因為舖入面冇乜人，得幾個閒逛嘅、一個
\VUgras{龍騎兵} \textsuperscript{(49)}，同一個\VUgras{收數員} \textsuperscript{(50)}——佢喺
\VUgras{總收銀處} \textsuperscript{(51)} 度遞緊張\VUgras{匯票。}

%%K t16-10-1 p
10. --- 我問阿爺，個\VUgras{女收銀}後面枱上面啲簿係做乜㗎；佢話我知嗰啲係
\VUgras{賬簿。}

%%K t16-tit-4 sub
\VUsaut{3.6mm}
\VUpk{10.2pt}{40}{圍住舖頭望。}
\VUsaut{2.8mm}

%%K t16-11-1 p
11. --- 離開玩具部之後，我哋去咗鞍具部望下啲\VUgras{馬鞍} \textsuperscript{(53)}、
\VUgras{馬鐙} \textsuperscript{(54)}、\VUgras{馬轡} \textsuperscript{(55)}、\VUgras{馬嚼}
\textsuperscript{(56)} 同啲\VUgras{馬鞭。}\VUgras{部長} \textsuperscript{(57)} 同我阿爺講緊
幾樣嘢嗰陣，我就行開去睇啲\VUgras{瓷器}同\VUgras{水晶} \textsuperscript{(58)} 嘅陳列
——嗰度有靚\VUgras{沙律盆，}同精緻嘅\VUgras{茶壺} \textsuperscript{(59)}。

%%K t16-12-1 p
12. --- 想上一樓，我哋要行過\VUgras{束腰} \textsuperscript{(60)} 嘅櫥窗後面，就喺
賣襪嗰檔隔籬——嗰度擺住好靚嘅\VUgras{長褲} \textsuperscript{(63)}。就喺附近，一個
\VUgras{女售貨員} \textsuperscript{(61)} 正畀個\VUgras{女顧客} \textsuperscript{(62)} 揀靚絲
\VUgras{綢} \textsuperscript{(64)}。

%%K t16-12-2 p
行樓梯上去嗰陣，我留意到成間舖用日本\VUgras{燈籠} \textsuperscript{(65)}、
\VUgras{陽傘} \textsuperscript{(66)} 同好大棵嘅\VUgras{棕櫚} \textsuperscript{(70)} 做裝飾。

%%K t16-13-1 p
13. --- 一樓冇乜好睇：淨係啲傢俬 (*)、啲\VUgras{夾萬} \textsuperscript{(67)}、啲
\VUgras{五桶櫃} \textsuperscript{(68)}，

%%K t16-13-1 p suite
同啲\VUgras{嬰兒車} \textsuperscript{(69)}。不過成間舖睇落去，就真係好\VUgras{好玩。}

%%K t16-14-1 p
14. --- 喺我哋腳下，我見到啲樂器：\VUgras{豎琴} \textsuperscript{(71)}、\VUgras{結他}
\textsuperscript{(72)}、\VUgras{曼陀林} \textsuperscript{(73)}、\VUgras{短號} \textsuperscript{(74)}、
\VUgras{單簧管} \textsuperscript{(75)}、小提琴連佢哋支\VUgras{弓} \textsuperscript{(76)}，連
\VUgras{大鼓} \textsuperscript{(77)} 都有。行遠少少，喺文具部，有個\VUgras{學生}
\textsuperscript{(78)} 同個\VUgras{男售貨員} \textsuperscript{(79)} 講緊本簿嘅價。佢哋隔籬，
我見到一本靚\VUgras{字母書} \textsuperscript{(80)}。

%%K t16-14-2 p
成間舖中間豎起一棵好高嘅\VUgras{聖誕樹} \textsuperscript{(81)}，掛滿蠟燭、細擺設同
各式各樣嘅玩具：\VUgras{木馬} \textsuperscript{(82)}、\VUgras{公仔} \textsuperscript{(83)}，如此類推。

%%K t16-14-3 p
\VUgras{香水部} \textsuperscript{(84)} 嗰啲\VUgras{牙膏}同各種\VUgras{香水}散出好香嘅味，
一路飄到我哋度。

%%K t16-tit-5 sub
\VUsaut{3.4mm}
\VUpk{10.2pt}{40}{我哋買咗啲嘢。}
\VUsaut{2.8mm}

%%K t16-15-1 p
15. --- 我阿爺開始覺得逛得太耐，我哋就行大樓梯落去。佢喺一幅
\VUgras{天文圖表} \textsuperscript{(85)} 前面停咗一陣，話畀我聽上面畫嘅係\VUgras{彗星}
\textsuperscript{(a)}、\VUgras{月食同日食} \textsuperscript{(b)}、標住\VUgras{四個正方位}
\textsuperscript{(c)} \VUgras{（北、南、東、西）}嘅風向玫瑰、兩個\VUgras{半球}
\textsuperscript{(d)} 嘅地圖、啲\VUgras{行星} \textsuperscript{(e)} 同個\VUgras{太陽系}
\textsuperscript{(f)}。呢啲嘢我一樣都唔明；行過嗰個着住鑲金線制服嘅
\VUgras{送貨夥計} \textsuperscript{(86)}，同埋我隔籬啲靚\VUgras{扇} \textsuperscript{(87)}——就喺
啲\VUgras{銀包} \textsuperscript{(88)} 同啲\VUgras{公事包} \textsuperscript{(89)} 隔籬——反而好睇好多。

%%K t16-16-1 p
16. --- 我又欣賞咗陳列上面啲\VUgras{羽毛} \textsuperscript{(90)}、啲\VUgras{皮帶扣}
\textsuperscript{(91)}，同啲靚靚哋插喺籮入面嘅\VUgras{假花} \textsuperscript{(94)}。啲
\VUgras{紫羅蘭、}

%%K t16-16-1 p suite
\VUgras{桂竹香、}\VUgras{風信子} \textsuperscript{(95)}、\VUgras{天竺葵} \textsuperscript{(96)}、
\VUgras{百合} \textsuperscript{(97)} 同\VUgras{金蓮花} \textsuperscript{(98)} 都做到十足十。

%%K t16-tit-6 sub
\VUsaut{4.4mm}
\VUpk{10.2pt}{40}{阿爺送我嘅禮物。}
\VUsaut{2.8mm}

%%K t16-17-1 p
17. --- 我求阿爺買一支\VUgras{牙刷} \textsuperscript{(92)} 畀我——嗰啲牙刷擺喺個盒度，
就喺啲\VUgras{髮夾} \textsuperscript{(93)} 隔籬。佢應承咗，我就自己攞去收銀處畀錢；
收銀處隔籬，人哋正執緊一大批\VUgras{貨} \textsuperscript{(100)}，準備寄畀一位
\VUgras{客} \textsuperscript{(99)}。

%%K t16-18-1 p
18. --- 突然之間，企喺啲藝術\VUgras{銅像} \textsuperscript{(101)} 前面嗰班人有少少騷動。
一個\VUgras{賊} \textsuperscript{(102)} 想扒人荷包，畀個\VUgras{巡查員} \textsuperscript{(103)}
當場捉住。人哋叫咗差人嚟\VUgras{拉}佢；我哋行出去嗰陣，個犯已經畀人押住入
差館。

\VUsaut{6.0mm}
\VUfilet{22mm}
